\section{Equational validity is coNP-complete}\label{sec:complexity}

The decision problem takes two ordinary tree terms, with any number of
variables, and asks whether they agree at every nonempty-set valuation.
Fixing $K$ fixes the algebra and its signature. In the uniform version,
role names have binary encodings and permutations are explicit lists of
their images. A succinct encoding of terms by shared recursive
descriptions is not assumed.

\begin{theorem}[Equational complexity]\label{gen:thm:complexity}
For every fixed $K\geq3$, equational validity is coNP-complete, already
in the signature with only variables and owned choices, without winner
constants or role switches. The uniform problem with $K$
as part of the input is coNP-complete as well.
\end{theorem}
\begin{proof}
For fixed $K$, a failed valuation is a polynomial certificate since the
algebra is finite and term evaluation is linear up to fixed constants.
For variable $K$, use the pair countervaluation theorem. A certificate
lists at most two role labels per input variable. Evaluate terms using
sparse sets: an owned node either collapses or unions its inputs, while
a permutation preserves cardinality. Each intermediate set has size at
most twice the number of variable leaves below it plus the number of
winner leaves. This is polynomial in tree-term size, independently of
the exponential size of the carrier and even when $K$ is binary encoded.
Evaluation and comparison are therefore polynomial. Validity is in coNP.

For hardness fix distinct roles $a,b$ and put $P(S)=[S=\{a\}]$.
On arbitrary nonempty inputs,
\[
 P(S\ow a T)=P(S)\lor P(T),\qquad
 P(S\ow b T)=P(S)\land P(T).
\]
The first equality is the absorbing-owner test. For the second, a
collapse to $\{b\}$ cannot pass $P$; otherwise a union is $\{a\}$ exactly
when both nonempty operands are $\{a\}$. Compile a constant-free monotone
Boolean formula $F$ to $\widehat F$ by replacing disjunction and
conjunction with these choices and leaving variables unchanged.
Then $P(\widehat F(v))=F(P\circ v)$ for every valuation, by induction.

The standard dual-rail reduction gives hardness directly. Given a CNF
formula, handle an empty clause by the fixed valid equation $x=x$,
and a CNF with no clauses by the fixed invalid equation $x=y$.
For the remaining $\varphi(z_1,\ldots,z_n)$, replace positive occurrences
of $z_i$ by $p_i$ and negative occurrences by $n_i$, producing a monotone
formula $\Phi(p,n)$. Put
\[
 A=\left(\bigwedge_i(p_i\lor n_i)\right)\land\Phi(p,n),
 \qquad B=\bigvee_i(p_i\land n_i).
\]
The equation $A\lor B=B$ fails exactly when $A=1$ and $B=0$.
Those conditions choose exactly one rail for every variable and satisfy
$\varphi$. Conversely any satisfying assignment provides such rails.
To reflect this Boolean equation as equality of winner sets, let $W$
be the union of all rail variables, expressed using a balanced fold of
$J$. Every constant-free term $u$ over those variables satisfies
$u(v)\subseteq W(v)$. Hence the constant-free gate
\[
 G_W(u)=u\ow a W
 =\begin{cases}
 \{a\},&u=\{a\}\text{ or }W=\{a\},\\
 W,&\text{otherwise}
 \end{cases}
\]
gives the identity
\[
 G_W(\widehat A\ow a\widehat B)=G_W(\widehat B).
\]
If $\varphi$ is unsatisfiable, the Boolean predicates agree on every
valuation. When $W=\{a\}$ both sides are $\{a\}$; otherwise the gate
maps equal predicates to the same winner set. If $\varphi$ is satisfiable,
assign each true rail $\{a\}$ and each false rail $\{b\}$.
Then $W=\{a,b\}$, $A=1$, and $B=0$, so the displayed identity fails.
These formulas use no Boolean constants in $A,B$.

Thus CNF unsatisfiability reduces to constant-free equation validity.
It is coNP-complete by the standard SAT theorem \cite{Cook1971}.
The balanced expansion of $W$ has $O(n^2)$ size: each $J$ uses its two
arguments twice. Formula sizes and their owned-choice translations are
polynomial, with no constants or switches. Hardness at $K=3$ also gives
uniform hardness.
\end{proof}

The gate uses the common input union as its false value, so it does not
require a named winner or $I$. The complexity method is standard. Unless $\mathrm P=\mathrm{NP}$,
the entire idempotent core has no polynomial-time canonicalization with
polynomial-time comparison. This does not preclude a finite equational
basis or an effective proof calculus.

\subsection{Hardness without owners}

The preceding reduction uses owned choices. Hardness also appears in
smaller signatures whose primitive operations are derived from them.
Use $A\del B=B$ when $B$ is a singleton, and $A\cup B$ otherwise;
its owned-choice definition and intrinsic characterization are given in
\cref{sec:reset}.
Here size is measured in the stated native signature; an expansion
into owned-choice syntax need not preserve that size.

Write $\operatorname{rc}(a_1,\ldots,a_n;b)$ for the right comb
$a_1\del(a_2\del\cdots\del(a_n\del b))$, and
$\operatorname{lc}(b;g_1,\ldots,g_m)$ for the left comb
$(\cdots((b\del g_1)\del g_2)\cdots)\del g_m$.
Both are $b$ when their lists are empty.

\begin{theorem}[Owner-free hardness]\label{thm:owner-free-hardness}
For every fixed $K\geq3$, equation validity for the signature
$\{\del,\cup\}$ on $\Pow(I)$ is coNP-complete. It uses no constants,
owner names, or internal role switches.
\end{theorem}
\begin{proof}
Membership in coNP follows either from a full valuation at fixed $K$,
or from the pair theorem and polynomial native evaluation. The pair
witness exists because both operations are owned-choice definable;
its existence does not require a small expansion of the input term.

For hardness, take a CNF $\varphi$ on $n\geq1$ Boolean variables.
Use literal variables $u_i,\bar u_i$ and three additional input variables
$t,z,c$. Put
\begin{align*}
 W&=\operatorname{rc}(u_1,\bar u_1,\ldots,u_n,\bar u_n,t;z),\\
 b_i&=(c\del(u_i\cup t))\del(\bar u_i\cup t),\\
 R&=\operatorname{rc}(b_1,\ldots,b_n;W),\\
 g_C&=t\cup\bigcup_{\ell\in C}\ell,\qquad
 A=\operatorname{lc}(c;(g_C)_{C\in\varphi}).
\end{align*}
The equation is
\begin{equation}\label{eq:union-reduction}
 A\del R=R.
\end{equation}
All displayed symbols are input variables or the two native operations.
The tree size is linear in the CNF description and $n$.

If $z$ is resolved, $W$, $R$, and both sides of
\eqref{eq:union-reduction} are $z$. Otherwise $W$ is the union of
$z,t$, and all literals; $R$ is ambiguous and contains $W$.
Every contribution other than $c$ is therefore already in $R$.
Outside $W$, a marker from $c$ survives in $A$ exactly when every
clause guard $g_C$ is ambiguous. It survives in $R$ exactly when
some pair $u_i\cup t,\bar u_i\cup t$ is simultaneously ambiguous.

Suppose the equation fails. Every clause guard is ambiguous, and no
pair is simultaneously ambiguous. Since $n\geq1$, $t$ must be a
singleton: an ambiguous $t$ would make every pair ambiguous.
For nonempty $X$, $X\cup t$ is resolved exactly when $X=t$.
Assign variable $i$ true iff $u_i\ne t$. If a positive literal differs
from $t$, it is true under this assignment. If a negative literal
$\bar u_i$ differs from $t$, consistency forces $u_i=t$, so that
negative literal is true. Every ambiguous clause guard has such a
literal, hence $\varphi$ is satisfied.

Conversely choose three distinct roles $b,d,e$ and a satisfying
assignment. Set $t=\{b\}$, each false literal to $\{b\}$, each true
literal to $\{d\}$, $z=\{b,d\}$, and $c=\{e\}$.
All clause guards are ambiguous, while every literal pair has a
resolved guard. Thus $e$ survives on the left and is absent on the
right. This proves the reduction. CNFs with no variables are decided
trivially and mapped to a fixed valid or invalid equation.
\end{proof}

\subsection{Hardness inside the atom-preserving fragment}

Put $H_s(x,y,z)=(y\star_s z)\del x$, the ternary witness used in
\cref{reset:thm:binary}. Each $H_s$ is idempotent and atom-preserving,
and reset is already definable by
\[
 x\del y=H_s(y,x,x).
\]
Thus the following result uses a signature of only two ternary operations.

\begin{theorem}[Two atom-preserving operations suffice]\label{thm:H-hardness}
For $K\geq3$ and distinct $a,b\in I$, equation validity for
$\{H_a,H_b\}$ on $\Pow(I)$ is coNP-complete, without constants.
\end{theorem}
\begin{proof}
The same finite-valuation argument gives membership in coNP.
For a CNF $\varphi$, use one input $y_i$ per Boolean variable, together
with $z,c$. Let $W=\operatorname{rc}(y_1,\ldots,y_n;z)$.
For a clause $C=(\ell_1,\ldots,\ell_m)$ define
\[
 b_C=H_{s_1}(z,y_{i_1},H_{s_2}(z,y_{i_2},\ldots
                  H_{s_m}(z,y_{i_m},c)\ldots)),
\]
where $s_j=a$ for a positive literal and $s_j=b$ for a negative one.
An empty clause gives $b_C=c$. Put
$R=\operatorname{rc}((b_C)_{C\in\varphi};W)$. The equation is
\begin{equation}\label{eq:H-reduction}
 c\del R=R.
\end{equation}
Expand each reset by $x\del y=H_a(y,x,x)$. In each right comb, the
continuation $y$ occurs once and only the small left operand is
repeated. The clause terms also use their continuation once. The
expanded tree therefore has polynomial, in fact linear, size in
$n$ and the total literal count.

When $z$ is resolved, both sides are $z$. Otherwise $W$ and $R$ are
ambiguous, and $W$ contains $z$ and every $y_i$. Outside $W$ the
exact propagation law is
\[
 H_s(z,y,u)\text{ contains a marker from }u
 \quad\Longleftrightarrow\quad
 u\text{ contains it and }y\ne\{s\}.
\]
This also holds when the marker is $s$: then $y\ne\{s\}$ follows
already from $y\subseteq W$. It follows by induction that a marker
from $c$ survives through $b_C$ iff all its literal tests
$y_{i_j}\ne\{s_j\}$ hold.

If $\varphi$ is unsatisfiable, read variable $i$ as true iff
$y_i=\{a\}$. Some clause is false. Its positive literals have
$y_i\ne\{a\}$, and its negative literals have $y_i=\{a\}\ne\{b\}$.
Consequently that branch retains every marker outside $W$ from $c$.
Since $W\subseteq R$, this gives $c\subseteq R$ and proves
\eqref{eq:H-reduction} on every valuation.

For a satisfying assignment, set $y_i=\{a\}$ for true and
$y_i=\{b\}$ for false, $z=\{a,b\}$, and $c=\{d\}$ for a third
role $d$. Every clause has a literal with $y_i=\{s\}$, removing
$d$ in its branch. Thus $d\notin R$, while $d\in c\del R$.
\end{proof}

These results do not exclude all tractable owner-sensitive fragments.
They do exclude polynomial-time equality for reset with union and for
the full idempotent atom-preserving clone unless $\mathrm P=\mathrm{NP}$.
Reset with one fixed $H_a$ is a separate unresolved case.

\subsection{The number of roles becomes observable again}

Let $\mathcal B_K=(\Pow(\{1,\ldots,K\});\del,\cup)$.
Unlike reset alone, these owner-free algebras have distinct equational
theories at every $K\geq3$.

\begin{theorem}[Strict dependence on the role count]\label{thm:colouring-theories}
For every $K\geq3$,
$\operatorname{Eq}(\mathcal B_{K+1})\subsetneq
 \operatorname{Eq}(\mathcal B_K)$.
\end{theorem}
\begin{proof}
For $n\geq3$ use variables $x_1,\ldots,x_n,z,c$. Define
\begin{align*}
 W_n&=\operatorname{rc}(x_1,\ldots,x_n;z),\\
 R_n&=\operatorname{rc}(c\del x_1,\ldots,c\del x_n;W_n),\\
 A_n&=\operatorname{lc}(c;(x_u\cup x_v)_{1\leq u<v\leq n}).
\end{align*}
Take the equation $A_n\del R_n=R_n$.
If $z$ is resolved, it holds immediately. Otherwise $W_n$ is ambiguous
and contains all the $x_i$. If any $x_i$ is ambiguous, $R_n$ also
contains $c$, so the equation holds. If all $x_i$ are singletons and
two agree, their edge guard is resolved and erases $c$ from $A_n$;
every remaining contribution lies in $W_n\subseteq R_n$.
The only remaining case has $n$ distinct singleton inputs. At $K=n$
they exhaust the carrier's roles, so $W_n=I$ and the equation holds.

At $K=n+1$ set $x_i=\{i\}$, $z=\{1,2\}$, and $c=\{n+1\}$.
Every edge guard is ambiguous; $A_n$ retains the fresh winner $n+1$,
while $R_n=W_n=\{1,\ldots,n\}$ omits it. The equation fails.
Finally, the inclusion of role sets induces a subalgebra embedding
$\mathcal B_K\hookrightarrow\mathcal B_{K+1}$: union and singleton
status are preserved. Equations valid in the latter hold in the
former, and the displayed equation makes this inclusion strict.
\end{proof}

The edge guards express proper colourings of a complete graph. A fresh
marker requires an unused role, which is precisely what fails at $K=n$.
Thus role dependence can return even when no operation names an owner.
